The ideal-gas theory begins with Bose's counting of light quanta
\cite{Bose1924} and Einstein's extension to a monatomic ideal gas,
including the condensation prediction in his second paper
\cite{Einstein1924IdealGas,Einstein1925IdealGas}.
Bogoliubov's treatment of weakly interacting bosons introduced an
approximate quasiparticle description based on macroscopic occupation of
one mode \cite{Bogolubov1947}.

For interacting particles, Penrose and Onsager formulated condensation
through a macroscopic eigenvalue of the one-particle density matrix and
applied this viewpoint to liquid helium \cite{PenroseOnsager1956}.
Yang developed the broader framework of off-diagonal long-range order,
covering Bose condensation and fermion pairing \cite{Yang1962ODLRO}.
These criteria identify the order to be proved; establishing it for a
specified interacting Hamiltonian remains a separate mathematical problem.

Rigorous lattice antecedents include Dyson, Lieb, and Simon's
low-temperature long-range order for the spin-$\tfrac12$ nearest-neighbor
quantum XY model on the simple cubic lattice in dimensions $d\ge3$
\cite{DysonLiebSimon1978}, and Kennedy, Lieb, and Shastry's ground-state
long-range order for all spins on hypercubic lattices in dimensions $d\ge2$
\cite{KennedyLiebShastry1988XY}.
For spin-$\tfrac12$, the spin--hard-core-boson correspondence relates
these results to condensation of half-filled lattice bosons.
Passing to a continuum Hamiltonian with a fixed interaction
and a prescribed density requires additional arguments.

The distinction between the thermodynamic and Gross--Pitaevskii limits is
essential. Lieb and Seiringer proved complete condensation in the
Gross--Pitaevskii limit, with the interaction scaled so that $Na$ remains
fixed, where $a$ is the scattering length
\cite{LiebSeiringer2002BEC}. Deuchert and Seiringer subsequently treated
positive-temperature homogeneous gases with scattering length
$a_N=a_vL/N$, where $a_v$ is the scattering length of the
unscaled potential, and controlled the one-particle density matrix of approximate
Gibbs minimizers \cite{DeuchertSeiringer2020GP}. These results allow the
interaction scale to vary with particle number. The problem considered
here instead fixes the potential and a positive density before taking
the volume to infinity.

The limits between Gross--Pitaevskii scaling and a fixed-density
thermodynamic limit have several forms. Brennecke, Brooks, Caraci, and
Oldenburg proved complete ground-state condensation in the
Gross--Pitaevskii regime and beyond
\cite{BrenneckeBrooksCaraciOldenburg2025}. Fournais treated low-energy
states on growing boxes in a joint dilute and increasing-volume limit
\cite[Theorem~1.2]{Fournais2021}. Chong, Liang, and Nam gave simpler
proofs of low-energy condensation using kinetic localization and
Poincar\'e-type inequalities \cite{ChongLiangNam2026}.
For the unit-box interaction $N^{2-2\kappa}v(N^{1-\kappa}x)$,
rescaling to a fixed potential gives density $N^{3\kappa-2}$;
the regimes $\kappa<2/3$ therefore couple dilution to increasing
particle number. Junge's low-temperature canonical Gibbs-state
result uses Neumann boxes whose size remains coupled to
the dilute limit \cite[Corollary~6]{Junge2026Propagation}.

The thermodynamic energy problem has a separate history. Dyson's
hard-sphere bounds and the leading dilute asymptotics of Lieb and
Yngvason established the first-order energy scale
\cite{Dyson1957,LiebYngvason1998}. The next correction was predicted
by Lee, Huang, and Yang's pseudopotential calculation
\cite{LeeHuangYang1957}. Fournais and Solovej proved the corresponding
lower bound for nonnegative radial finite-range potentials, including hard cores
\cite[Theorems~1.1 and~1.3]{FournaisSolovej2023}; recent upper bounds
attain the same coefficient for hard spheres and for general
nonnegative even finite-range potentials
\cite{BastiEtAl2026HardSpheres,BastiEtAl2026GeneralPotentials}.
These energy results, and low-temperature free-energy estimates for
strong potentials \cite{FournaisEtAl2026StrongFreeEnergy}, do not by
themselves determine occupation of a single orbital: the one-particle
kinetic gap tends to zero as the volume grows.

A distinct thermodynamic-limit result appears in S\"ut\H{o}'s two-step preprint
\cite{Suto2023TwoStep}. Its Theorem~1.2 assumes nonnegativity both of the
potential and of its Fourier transform, together with integrability,
a Fourier-moment condition, and spatial decay, and states condensation
above a temperature-dependent density threshold. Fourier positivity is
an additional restriction: for example, a nonnegative spherical square
barrier has a Fourier transform that changes sign. That result therefore
does not cover all bounded nonnegative radial finite-range potentials,
and its fixed-temperature threshold does not alone provide the joint
low-density and low-temperature uniformity sought here.

Galanda and Pinamonti construct algebraic equilibrium states with a
condensate background under quantitative convergence conditions for a
spatially localized interaction \cite[Theorem~6.8]{GalandaPinamonti2025Equilibrium}.
Their subsequent analysis treats renormalized correlation functions in
Gross--Pitaevskii scaling, with a zero-temperature limit in its final
convergence theorem \cite[Theorem~7.3]{GalandaPinamonti2026GP}.
These constructions use a different state formulation and coupled
interaction, regulator, and temperature limits from the fixed-potential
canonical thermodynamic limit considered here.

The trajectory formulation descends from Feynman's treatment of the
Bose gas \cite{Feynman1953AtomicLambda}; for the canonical
Feynman--Kac representation and its symmetrization and one-particle
normalizations, see \cite[Appendix~A]{Ueltschi2006FeynmanCycles}.
The random-path comparison has a different ancestry.
Benjamini, Pemantle, and Peres constructed paths with exponentially
decaying intersection tails \cite{BenjaminiPemantlePeres1998}.
H\"aggstr\"om and Mossel developed a hierarchical random-drift
construction with low predictability; its independent block variables
are a close antecedent of the dyadic velocities used below
\cite[Section~3, Proposition~3.1]{HaggstromMossel1998}.
Abbe, Massouli\'e, Montanari, Sly, and Srivastava used finite-endpoint
path averages and overlap estimates for group synchronization
\cite{AbbeMassoulieMontanariSlySrivastava2018}; Garban and Spencer used
this approach for classical disordered-spin symmetry breaking along the Nishimori line
\cite{GarbanSpencer2023}. Here the finite cell paths, their thickened
intersection estimates, and the continuum bridge comparison require
separate proofs.
